\section{ЛР 1. База знаний Prolog}
\subsection{Предметная область и словарь}
Объекты представлены атомами в snake\_case: например, \code{portal_2},
\code{valve}. Переменные начинаются с заглавной буквы. Факт утверждает отношение,
а правило задаёт условия его доказательства. Запятая означает конъюнкцию,
точка с запятой --- дизъюнкцию. Отрицание \code{\+ Goal} успешно, когда
интерпретатор не смог доказать цель Goal.

\begin{table}[H]\centering
\caption{Факты исходной базы}\begin{tabular}{lr}\toprule
Предикат & Число фактов\\\midrule
\code{game/1} & 12\\
\code{single_player/1} & 11\\
\code{multiplayer/1} & 6\\
\code{story_driven/1} & 7\\
\code{difficult/1} & 5\\
\code{indie/1} & 4\\
\code{relaxing/1} & 3\\
\code{competitive/1} & 2\\
\code{free_to_play/1} & 1\\\midrule
Всего унарных & 51\\
\code{developed_by/2} & 10\\
\code{has_genre/2} & 5\\\bottomrule
\end{tabular}\end{table}

\subsection{Правила и порядок вычисления}
В файле games.pl факты сгруппированы по признакам и отношениям,
далее записаны правила. Группы фактов и основные правила снабжены комментариями.
Примеры запросов вынесены в queries.md и run\_examples.pl,
автоматические проверки --- в test\_games.pl; загрузка БЗ сама по себе
не запускает демонстрации или тесты.

Семь правил описывают сюжетные одиночные игры, дружелюбные многопользовательские,
напряжённые, сюжетные независимые, связь игры со студией, спокойные одиночные
и подходящие новичку игры. Пример композиции условий:
\begin{lstlisting}
recommended_for_beginner(Game) :-
    game(Game), single_player(Game),
    \+ difficult(Game), \+ competitive(Game).
\end{lstlisting}
Сначала выбирается конкретная игра, затем проверяются отрицания. Такая последовательность
существенна: проверка отрицания со свободной переменной спрашивала бы об отсутствии
любой сложной игры, а не о свойствах кандидата.
В \code{intense_game/1} игра Civilization VI удовлетворяет двум ветвям ИЛИ;
\code{setof/3} устраняет повторения в итоговом наборе.

\subsection{Отсечение и полнота}
Предикат \code{!} отсекает альтернативы, возникшие до него в пределах текущего вызова.
В запросе \code{story_driven(G), !, \+ developed_by(G,cd_projekt_red)}
первым выбирается The Witcher 3. Отрицание не выполняется, а возвращаться к
следующим сюжетным играм уже нельзя, поэтому ответ --- false.
Без отсечения находятся Portal 2, Hades, Hollow Knight, Celeste и Baldur's Gate 3.
Здесь cut меняет множество решений и не подходит для получения всех рекомендаций.

\subsection{Фактические ответы интерпретатора}
Ниже приведены десять запросов, выполненных скриптом сбора результатов.
\code{findall/3} собирает ответы в порядке поиска; \code{setof/3} возвращает
отсортированный набор. Списки совпадают с отдельными интерактивными решениями.
\par\medskip
\noindent\begin{minipage}{\linewidth}
\textbf{Запрос 1}
\begin{lstlisting}
?- game(minecraft).
true
\end{lstlisting}
\end{minipage}\par\smallskip

\noindent\begin{minipage}{\linewidth}
\textbf{Запрос 2}
\begin{lstlisting}
?- developed_by(portal_2,valve).
true
\end{lstlisting}
\end{minipage}\par\smallskip

\noindent\begin{minipage}{\linewidth}
\textbf{Запрос 3}
\begin{lstlisting}
?- findall(G,developed_by(G,valve),L).
[portal_2, counter_strike_2]
\end{lstlisting}
\end{minipage}\par\smallskip

\noindent\begin{minipage}{\linewidth}
\textbf{Запрос 4}
\begin{lstlisting}
?- findall(G,(single_player(G),story_driven(G)),L).
[the_witcher_3, cyberpunk_2077, portal_2, hades, hollow_knight,
celeste, baldurs_gate_3]
\end{lstlisting}
\end{minipage}\par\smallskip

\noindent\begin{minipage}{\linewidth}
\textbf{Запрос 5}
\begin{lstlisting}
?- findall(G,(game(G),(relaxing(G);free_to_play(G)),\+ competitive(G)),L).
[minecraft, stardew_valley]
\end{lstlisting}
\end{minipage}\par\smallskip

\noindent\begin{minipage}{\linewidth}
\textbf{Запрос 6}
\begin{lstlisting}
?- findall(G,friendly_multiplayer(G),L).
[portal_2, minecraft, stardew_valley, baldurs_gate_3]
\end{lstlisting}
\end{minipage}\par\smallskip

\noindent\begin{minipage}{\linewidth}
\textbf{Запрос 7}
\begin{lstlisting}
?- setof(G,intense_game(G),L).
[celeste, civilization_vi, counter_strike_2, doom_eternal, hades,
hollow_knight]
\end{lstlisting}
\end{minipage}\par\smallskip

\noindent\begin{minipage}{\linewidth}
\textbf{Запрос 8}
\begin{lstlisting}
?- findall(G,indie_gem(G),L).
[hades, hollow_knight, celeste]
\end{lstlisting}
\end{minipage}\par\smallskip

\noindent\begin{minipage}{\linewidth}
\textbf{Запрос 9}
\begin{lstlisting}
?- findall(G,recommended_for_beginner(G),L).
[the_witcher_3, cyberpunk_2077, portal_2, minecraft, stardew_valley,
baldurs_gate_3]
\end{lstlisting}
\end{minipage}\par\smallskip

\noindent\begin{minipage}{\linewidth}
\textbf{Запрос 10}
\begin{lstlisting}
?- findall(G,relaxing_solo_game(G),L).
[minecraft, stardew_valley, civilization_vi]
\end{lstlisting}
\end{minipage}\par\smallskip

